\section{Localisation uniqueness on the selected integral Hecke module}

% BEGIN L84C1
\paragraph{L84C1. The character and row comparisons on an original member.}
Fix any selected member $D=w.D(Z)$, rows $t$, norm choice $n$ and
extension witness $c$ in the original member quantifiers.
The record $c.a.nt$ presents this very $D$ as the localisation of its
row and character data. Its Weil--Deligne identification at a place
$v$ preserves dimension, zero, sum, tensoring by a character,
real unramified twist, unramifiedness and Frobenius eigenvalues.
In the identified category a line has zero nilpotent operator.
Tensoring by a line distributes over a direct sum, commutes with
tensoring by the norm to a real power on either factor, and
successive character tensors compose by multiplication.
These isomorphisms intertwine both the Weil actions and the
nilpotent operators. The unramified line with Frobenius eigenvalue
one is the trivial line, so tensoring by it is the identity.
The bijection on isomorphism classes transports all five assertions
back to the member. These are the five laws of
\(\mathrm{RowTie.CharLaws}\).

At every split $v$ the two specialisations of a common row $\sigma$
at any $i,i'$ have local parameter equal to the one underlying
row parameter $\mathrm{locR}(\sigma,v)$. If $j_i=j_{i'}=j$ and
$w_i=w_{i'}=w$, the two tuple inputs $\mu_i,\mu_{i'}$ both equal
$\mu(w)$. In the selected character frame the two automorphic
characters $\mu_{C,i},\mu_{C,i'}$ both map to the $(j-1)$-indexed
entry of this tuple. Here $1\le j\le m$ makes this a common
zero-based entry. The frame identification is injective, so the
characters themselves agree. These comparisons prove the three
conjuncts of the built laws on the specified member.
% END L84C1

% BEGIN L84C2
\paragraph{L84C2. Comparing two membership polynomials before reduction.}
Suppose $i,i'$ have this same $(j,w)$ and $\sigma$ satisfies
$\mathrm{AtM}$ at both indices. Choose the two lists $\Sigma,\Sigma'$
and the two points $x,x'$ in their respective witnesses. At a common
split test place $v$, the lists separately have parameter
$\mathrm{locW}(i,\sigma,v)$ and $\mathrm{locW}(i',\sigma,v)$.
They are equal by L84C1. Appending $(1,\mu_{C,i},t)$ and
$(1,\mu_{C,i'},t)$ adds the same direct sum on both sides.
Indeed the term for a part $(a,\pi,d)$ is the sum of the local
parameters for $\pi$ with real twists $(d-1)/2-k$, $0\le k<d$.
Twisting the augmented list by $\eta_T^{-1}$ thus tensors the
entire sum by its local character. To see this including the empty
list, the zero local parameter tensored with this unramified
line still has dimension zero and empty Frobenius multiset;
the unramified parameter with that multiset is zero. Each added
summand uses the direct-sum and real-twist identities of L84C1
and compatibility of the local character with the global twist.
Finally apply the same real twist $(1-N)/2$, $N=n_V$, to both sides.
The resulting Galois polynomials over $\mathbb C$ agree.

Let $O_i,O_{i'}$ be the actual coefficient rings. The selected
localisation data identify both with $O_c$ for $(j,w)$, by
$\phi_i,\phi_{i'}$. The integral complex embeddings of the member
are the composites of these identifications with the common
embedding of $O_c$. Indeed the two given fraction-field embeddings
into the local algebraic closure and its complex embedding are
exactly the specialisations in the chosen data. The embedding on
$O_i$ is injective, as the composite of its map into its fraction
field and field embeddings. Consequently the polynomials $P_{i,v}$
and $P_{i',v}$ supplied by $\mathrm{AtM}$ have identical images in
$O_c[X]$. This is coefficientwise equality, before any use of the
Hecke congruence.
% END L84C2

% BEGIN L84C3
\paragraph{L84C3. The residual maximal ideal and the original localisation index.}
Fix the common $(j,w)$ of L84C2 and its selected Hecke algebra
$\mathcal H=O_c[H_{v,k}]$ inside the endomorphisms of the common
coefficient space. Its generators run over $\mathrm{GenAt}$.
At both indices they are split tests away from the residue
characteristic, the designated bad places and the common auxiliary
set $Q_c$. This uses $Q_i=Q_c=Q_{i'}$ and the equality of residue
characteristics forced by the common $w$. Under the selected
localisation map the scalar extension of the integral module
at $i$ injects into the common space with image the simultaneous
part at its maximal ideal $\mathfrak m_i$. Its point $x$
corresponds to an $O_c$-algebra eigencharacter
$\psi_x:\mathcal H\to O_c$ with nonzero eigenspace,
$\ker(\overline{\psi}_x)=\mathfrak m_i$, and
$\psi_x(H_{v,k})=\phi_i(h_{i,v,k}(x))$. The same description
applies to $i',x'$. All these identifications belong to the
selected normalised spectrum.

For every generator, compare the coefficient of $X^{N-k}$ in
the polynomials of L84C2. By the two congruences in $\mathrm{AtM}$,
after reduction in the common residue field it is on each side
the corresponding Hecke coefficient
\[
 (-1)^k q_v^{k(k-1)/2}\,\overline\psi_x(H_{v,k}),
 \qquad 1\le k\le N ,
\]
with $x'$ on the other side. This follows by expanding the monic
Hecke polynomial, since exactly the term with index $k$ has that
degree. Both coefficient isomorphisms respect the maximal ideals.
The scalar in front is a unit since $v\nmid\ell_i$. Thus the two
reduced eigencharacters agree on every generator and on $O_c$.
Every element of $\mathcal H$ is obtained by finite sums and
products of these elements, so the maps agree on $\mathcal H$.
Their kernels are equal. The map $i\mapsto\mathfrak m_i$ is injective
on the specified fibre of the selected normalised spectrum.
Therefore $i=i'$.

In the abstract Hecke-tied formulation the same argument uses the
given coefficient isomorphism compatible with complex embeddings.
If $i\ne i'$, its separation clause furnishes a common generator
whose two integral eigenvalues differ by a unit. L84C2's polynomial
comparison (using that coefficient isomorphism) and the unit-scalar
calculation above put the same difference in the maximal ideal.
This proves the abstract implication as well. On the original
member the internal spectrum proves the separation directly, since
different residual kernels must already differ on a generator.
In particular the result applies for every twisted $i$, every
$i'$ with the same $(j,w)$, and every good row $\sigma$ occurring
in $\mathrm{AtM}$ at both. This supplies the field $c2$ of
$\mathrm{GaloisPreconditions}$ on the same selected member.
% END L84C3
